% !PRINTED 482
\weberpage{482}

Ein System von drei Functionen, die sich ternär substituiren, kann man auf folgende Weise bilden\footnote{F.\ \gesperrt{Klein}, \glqq Ueber die Auflösung gewisser Gleichungen vom 7\textsuperscript{ten} und 8\textsuperscript{ten} Grade\grqq. Mathem. Annalen, Bd. XV (1879).}.

Wir führen zunächst ein System von Hülfsvariablen $x_1, x_2, x_3$ ein, die wir den Substitutionen der Gruppe $G_{168}$ unterwerfen, und daraus bilden wir die zur Gruppe $G_{24}$ gehörige Function $z_0$ mit ihren conjugirten $z_r$ [\S.\,125, (7)]:
\weq{2}{z_r = \varepsilon^{2r}x_1^2 + \varepsilon^{4r}x_2^2 + \varepsilon^r x_3^2 - \frac{1-\sqrt{-7}}{2}(\varepsilon^{-r}x_2x_3 + \varepsilon^{-2r}x_3x_1 + \varepsilon^{-4r}x_1x_2).}

Hierzu nehmen wir nun eine Function $\eta_0$ der Wurzeln $\lambda_r$ unserer Gleichung 7\textsuperscript{ten} Grades, die zu der Gruppe $P_{24}$ gehört, z.\,B. eine rationale Function von $\lambda_0$, und die conjugirten Werthe $\eta_0, \eta_1, \ldots, \eta_6$, und bilden die Summe
\weq{3}{\psi = \eta_0 z_0 + \eta_1 z_1 + \eta_2 z_2 + \eta_3 z_3 + \eta_4 z_4 + \eta_5 z_5 + \eta_6 z_6,}
die eine quadratische Function der $x$ ist, deren Coëfficienten von den Wurzeln $\lambda_r$ abhängen.

\gesperrt{Diese Function $\psi$ ändert sich nicht, wenn die Variablen $(x)$ einer Substitution der Gruppe $G_{168}$ und die Wurzeln $\lambda_r$ gleichzeitig der entsprechenden Permutation aus $P_{168}$ unterworfen werden.} Denn durch diese gleichzeitige Operation werden in der Summe (2) nur die Summanden unter einander vertauscht, also die Summe selbst nicht geändert.

Wir können daher $\psi$ als \gesperrt{simultane Invariante} der Gruppen $G_{168}$ und $P_{168}$ bezeichnen.

Ordnen wir die Function $\psi$ nach den Variablen $x_1, x_2, x_3$, so ergiebt sich
\weq{4}{\psi = p_1x_1^2 + p_2x_2^2 + p_3x_3^2 + 2q_1x_2x_3 + 2q_2x_3x_1 + 2q_3x_1x_2,}
worin zur Abkürzung
\weq{5}{\begin{aligned}
p_1 &= \sum^r \varepsilon^{2r}\eta_r, & q_1 &= -\frac{1-\sqrt{-7}}{4}\sum^r \varepsilon^{-r}\eta_r\\
p_2 &= \sum^r \varepsilon^{4r}\eta_r, & q_2 &= -\frac{1-\sqrt{-7}}{4}\sum^r \varepsilon^{-2r}\eta_r\\
p_3 &= \sum^r \varepsilon^{r}\eta_r, & q_3 &= -\frac{1-\sqrt{-7}}{4}\sum^r \varepsilon^{-4r}\eta_r
\end{aligned}}
